%
%  chapter1.tex — Introduction
%
\chapter{Introduction}
\label{cha:introduction}

\section{Context and Motivation}
\label{sec:context}

The world simultaneously faces an education crisis and an engagement crisis. Sustainability topics — climate action, responsible consumption, digital citizenship — are recognized as urgent priorities by the United Nations through the 2030 Agenda for Sustainable Development Goals (SDGs)~\cite{unsdgs2015}. Yet traditional educational methods fail to convert knowledge into lasting behavioral change. Passive lectures, static documents, and one-size-fits-all curricula produce low retention and even lower motivation in learners who are bombarded with information but rarely empowered to act on it.

Parallel to this, digital literacy has become a fundamental skill in a world where misinformation spreads rapidly and critical thinking about technology is essential. Educators need accessible tools to build and deliver quality learning experiences without requiring deep technical expertise. Learners need experiences that fit their pace, respect their time, and reward consistent effort.

These two challenges share a common root: the gap between content and engagement. Existing platforms either focus on content delivery without personalization, or on engagement mechanics without educational rigor. Play4Change is designed to close that gap.

\section{Problem Statement}
\label{sec:problem}

Current e-learning platforms suffer from well-documented limitations:
\begin{itemize}
    \item \textbf{Low engagement}: Completion rates for Massive Open Online Courses (MOOCs) rarely exceed 10\%~\cite{jordan2014initial}.
    \item \textbf{Lack of personalization}: Courses are delivered at a fixed pace regardless of the learner's prior knowledge or performance.
    \item \textbf{Content bottleneck}: Educators must manually author every quiz, challenge, and task, limiting scalability.
    \item \textbf{Weak feedback loops}: Learners receive little indication of their progress relative to their goals or peers.
    \item \textbf{High infrastructure cost}: Many platforms require expensive cloud resources that scale poorly with the user base.
\end{itemize}

From a systems engineering perspective, there is also a lack of rigorously architected platforms in the educational technology space. Most edtech solutions are built as monolithic applications with tight coupling, making them difficult to maintain, test, and scale.

\section{Objectives}
\label{sec:objectives}

Play4Change addresses these problems by pursuing the following primary objectives:

\begin{enumerate}
    \item Design and implement a \textbf{cross-platform mobile application} (iOS and Android) using Kotlin Multiplatform (KMP) with Decompose for lifecycle-aware navigation.
    \item Build a \textbf{scalable, domain-driven backend} capable of delivering personalized daily challenges at predictable cost.
    \item Integrate \textbf{generative AI} (Mistral LLM) with Retrieval-Augmented Generation (RAG) and vector search for automated, high-quality content creation from educator-supplied materials.
    \item Apply \textbf{clean architecture, hexagonal architecture, and Domain-Driven Design} (DDD) principles consistently across all system layers.
    \item Ensure the system prioritizes \textbf{security} (passwordless authentication, minimal sensitive data retention), \textbf{accessibility} (caching, low-bandwidth friendliness), and \textbf{future scalability} (Kubernetes-ready infrastructure, CDN-friendly design).
    \item Provide a \textbf{curated topic ecosystem} aligned with UN Sustainable Development Goals and digital literacy curricula.
\end{enumerate}

\section{Supervisors}
\label{sec:supervisors}

This project is supervised by:
\begin{itemize}
    \item \textbf{Prof.\ Nuno Datia} — ISEL, Engenharia Informática
    \item \textbf{Prof.\ António Serrador} — ISEL, Engenharia Informática
    \item \textbf{Prof.\ Michel Vorenhout} — Ureka / Sustainability domain co-supervisor
\end{itemize}

\section{Report Structure}
\label{sec:structure}

The remainder of this report is organized as follows:

\begin{description}
    \item[Chapter~\ref{cha:background}] \textbf{Background and State of the Art} — reviews the relevant literature on e-learning, gamification, adaptive learning, AI-generated education, sustainability curricula, and the architectural patterns employed in the system.
    \item[Chapter~\ref{cha:platform}] \textbf{Play4Change Platform} — describes the product vision, core features, user roles, and the four-step learning journey.
    \item[Chapter~\ref{cha:architecture}] \textbf{System Architecture} — presents the technical architecture, technology stack, component breakdown, and the engineering principles that govern design decisions.
    \item[Chapter~\ref{cha:implementation}] \textbf{Implementation} — details the current development state, testing strategy, CI/CD pipeline, and known challenges.
    \item[Chapter~\ref{cha:conclusion}] \textbf{Conclusion and Future Work} — summarizes the contributions, reflects on the project status, and outlines planned future work.
\end{description}
